\documentclass[11pt,reqno]{amsart}
\usepackage[T1]{fontenc}
\usepackage{lmodern}
\usepackage[margin=2.6cm]{geometry}
\usepackage{amsmath,amssymb,amsthm,mathtools,booktabs,array,enumitem,microtype}
\usepackage[colorlinks=true,linkcolor=blue,citecolor=blue,urlcolor=blue]{hyperref}
\setlength{\emergencystretch}{3em}
\allowdisplaybreaks

\newtheorem{theorem}{Theorem}[section]
\newtheorem{proposition}[theorem]{Proposition}
\newtheorem{lemma}[theorem]{Lemma}
\newtheorem{corollary}[theorem]{Corollary}
\theoremstyle{definition}
\newtheorem{definition}[theorem]{Definition}
\newtheorem{hypothesis}[theorem]{Hypothesis}
\theoremstyle{remark}
\newtheorem{remark}[theorem]{Remark}
\numberwithin{equation}{section}

\newcommand{\Z}{\mathbb Z}
\newcommand{\N}{\mathbb N}
\newcommand{\Q}{\mathbb Q}
\newcommand{\C}{\mathbb C}
\newcommand{\PP}{\mathbb P}
\newcommand{\A}{\mathbb A}
\newcommand{\ub}{\bar u}
\newcommand{\vb}{\bar v}
\newcommand{\Bred}{B_{\mathrm{red}}}
\newcommand{\Res}{\operatorname{Res}}
\newcommand{\lst}{\ell^{*}}
\newcommand{\tst}{t^{*}}
\newcommand{\Ncen}{N_k^{\mathrm{cen}}}

\title{Counting powerful products of consecutive integers}
\author{Danny Plymesser}
\date{Version 1.0.0, 23 October 2026}

\begin{document}

\begin{abstract}
We show that the number of $n\le x$ for which $n$ and $n+1$ are both powerful
is $O_\varepsilon(x^{1/4+\varepsilon})$, improving the exponent $29/100$ of
Reuss. For each fixed $k\ge 3$ we show that the number of $n\le x$ for which
$(n+1)(n+2)\cdots(n+k)$ is powerful is $O_{k,\varepsilon}(x^{77/320+\varepsilon})$,
and we prove the same exponent uniformly over all widths $k\ge3$. The new
ingredient is a per-interval count in Reuss's form of Heath-Brown's approximate
determinant method: in lattice-reduced coordinates, the solutions on an
auxiliary curve of bidegree $(\alpha,\beta)$ become integral points of height
at most $L_1^3T_1^2$ on an affine curve of degree $2\alpha+3\beta$, to which
the uniform bound of Castryck, Cluckers, Dittmann and Nguyen applies.
The proofs are formalized in Lean~4 with Mathlib and checked by Lean's kernel
and an independent kernel; the formal statements take six theorems from the
literature as explicit hypotheses. The results and exposition were produced and reviewed
using AI systems; no human mathematical review has been performed.
Section~\ref{sec:status} describes the verification and its scope.
\end{abstract}

\maketitle

\section{Introduction}\label{sec:intro}

A positive integer $m$ is \emph{powerful} (or square-full) if $p^2\mid m$ for
every prime $p\mid m$. Every powerful number can be written uniquely as $d^3u^2$
with $d$ squarefree, so there are $\asymp x^{1/2}$ powerful numbers up to $x$.

\subsection{Consecutive powerful numbers}
Pairs of consecutive powerful numbers exist in abundance: every solution of
$X^2-8Y^2=1$ gives the pair $8Y^2,\,X^2$, and this Pell family alone
contributes $\asymp\log x$ pairs up to $x$; see \cite{Golomb,Walker} for other
families. Erd\H{o}s conjectured that there are no three
consecutive powerful numbers \cite{Erdos76,ErdosManitoba}\cite[Problem \#364]{Bloom}, and it has been
conjectured that the number of consecutive pairs up to $x$ is
$(\log x)^{O(1)}$; see \cite{Erdos76}, \cite[p.~68]{ErdosGraham}, and
\cite[(1.15)]{Tao}, \cite[Problem \#365]{Bloom}. Erd\H{o}s also asked
whether a product of $k\ge3$ consecutive integers can be powerful
\cite[Problem \#137]{Bloom}; by the Erd\H{o}s--Selfridge theorem such a product is
never a perfect power \cite{ErdosSelfridge}. Under the $abc$ conjecture the
number of consecutive pairs is $x^{o(1)}$ and consecutive triples are finite in
number; see \cite[(1.16) and the following discussion]{Tao}.

Unconditionally, the trivial bound for
\[
N(x):=\#\{n\le x:\ n\text{ and }n+1\text{ are powerful}\}
\]
is $O(x^{1/2})$. Chan \cite{Chan2012} proved that the number of $n\le x$ with
$n$ and $n+l$ both powerful is $\ll d_3(l)x^{2/5}(\log x)^2$ uniformly in
$l\ge1$ \cite[Theorem~3]{Chan2012}, and $\ll_l x^{7/19}\log x$ for fixed $l$
\cite[Theorem~4]{Chan2012}. Akta\c{s} and Murty \cite{AktasMurty} gave
an $O(x^{2/5})$ bound for $N(x)$. Tao records
$x^{2/5}$ as the best known bound \cite[(1.16)]{Tao} and proves an estimate of
the same strength, $x^{2/5+o(1)}$, for the non-powerful integers lying in
intervals of length at least two with powerful product
\cite[Theorem~1.8]{Tao}. Reuss \cite[Theorem~4]{Reuss} proved
$N(x)\ll_\varepsilon x^{29/100+\varepsilon}$ with his extension of
Heath-Brown's approximate determinant method \cite{HBsqfree,HBAnnals}; the
exponent $29/100$ is also recorded in \cite{WongcharoenbhornMeemark}. We are
not aware of a published improvement of $29/100$.

\subsection{Results}
\begin{theorem}\label{thm:k2}
For every $\varepsilon>0$,
\[
\#\{n\le x:\ n\text{ and }n+1\text{ are powerful}\}\ll_\varepsilon x^{1/4+\varepsilon}.
\]
\end{theorem}

The exponent $1/4$ is the exact limit of our method for $k=2$
(Remark~\ref{rem:quarter}). For $k\ge2$ and $x\ge1$ put
\[
N_k(x):=\#\{0\le n\le x:\ (n+1)(n+2)\cdots(n+k)\text{ is powerful}\}.
\]
For $k\ge3$ the product can be powerful without any single factor being
powerful, because primes $p<k$ may divide two factors. Every $n$ counted by
$N_k(x)$ has, for $1\le i\le k$, a decomposition $n+i=f_id_i^3u_i^2$, where
$d_i$ is the product of the primes $p\ge k$ that divide $n+i$ to an odd power
(Lemma~\ref{lem:decomp}). Define
\begin{multline*}
\Ncen(x):=\#\{n\le x:\ (n+1)\cdots(n+k)\text{ is powerful, and}\\
(n+3)^{19/100}<d_i<(n+3)^{21/100}\ \text{for } i=1,2,3\}.
\end{multline*}
For $k=3$ this is the ``central box'' of the three-term problem, the region in
which elementary pairwise estimates are weakest.

\begin{theorem}\label{thm:k3}
Let $k\ge3$ be fixed. For every $\varepsilon>0$,
\[
N_k(x)\ll_{k,\varepsilon}x^{77/320+\varepsilon}
\qquad\text{and}\qquad
\Ncen(x)\ll_{k,\varepsilon}x^{38059/160000+\varepsilon}.
\]
\end{theorem}

Here $77/320=0.240625$ and $38059/160000=0.23786875$. We know of no bound for
$N_k(x)$, $k\ge3$, stated in the literature, but two bounds follow at once from
the literature: Chan's fixed-shift bound \cite[Theorem~4]{Chan2012} applied
to the pair equation of Lemma~\ref{lem:pair} gives $N_k(x)\ll_kx^{7/19}\log x$,
and \cite[Theorem~1.8 and (1.16)]{Tao} give $N_k(x)\ll x^{2/5+o(1)}$. Both
deductions are ours. Theorem~\ref{thm:k3} improves both exponents.

\begin{theorem}[Uniformly over all widths]\label{thm:uniform}
For every $\varepsilon>0$,
\[
 \#\{0\le n\le x:\exists k\ge3,\ P(n,k)\text{ powerful}\}
 \ll_\varepsilon x^{77/320+\varepsilon}.
\]
With $k\ge2$ in place of $k\ge3$, the bound is
$O_\varepsilon(x^{1/4+\varepsilon})$.
The number of non-powerful positive integers at most $x$ which belong to
an interval of consecutive positive integers with powerful product is
$O_\varepsilon(x^{77/320+\varepsilon})$.
The implied constants are not explicit: they depend on the threshold in
Theorem~\ref{thm:tao}.
\end{theorem}

Here $P(n,k)=(n+1)\cdots(n+k)$. The first bound counts the union of starts
over all widths, with one constant independent of $k$. The last assertion
bounds the quantity in Tao's Theorem~1.8.

\subsection{Method}
Write $n+i=f_id_i^3u_i^2$ and $n+j=f_jd_j^3u_j^2$. After absorbing the bounded
coefficients $f_i,f_j$, each pair of factors gives a solution of
\begin{equation}\label{eq:pair}
e^3v^2-d^3u^2=h,
\end{equation}
with $h\ne0$ bounded in terms of $k$ (Lemma~\ref{lem:pair}). Reuss
\cite[\S2]{Reuss} splits the range of $\rho=d/e$ into $\asymp M$ intervals
$I$, with $\log M=\tfrac98(1+\delta)\log(DE)\log(UV)/\log x$, and shows by the
determinant method that on each interval all solutions lie on one auxiliary
curve of bounded degree (Lemma~\ref{thm:lemma7}). He then passes to
lattice-reduced coordinates $(\lambda,\tau)\in\Z^2\times\Z^2$ with
$|\lambda_i|\le L_i$, $|\tau_i|\le T_i$, and counts the points of the
auxiliary curve using one factor at a time, which gives
$\ll x^{\varepsilon}\min(L_1,T_1)$ solutions per interval.

Our new step (Lemma~\ref{lem:P}) also uses the equation \eqref{eq:pair} itself.
If $\Gamma\subset\PP^1\times\PP^1$ is a component of bidegree $(\alpha,\beta)$ of
the auxiliary curve, the twelve monomials of bidegree $(3,2)$ in
$(\lambda,\tau)$ map the solutions on $\Gamma$, at most two-to-one, to
integral points of height at most $L_1^3T_1^2$ on the affine curve obtained by
slicing the cone over the Segre--Veronese image of $\Gamma$ with the
hyperplane $\{\ell=h\}$, where $\ell$ represents the left-hand side of
\eqref{eq:pair}. This curve has degree $2\alpha+3\beta$, and the bound of
Castryck, Cluckers, Dittmann and Nguyen \cite{CCDN} gives
$\ll(L_1^3T_1^2)^{1/(2\alpha+3\beta)}\log(2+L_1T_1)$ solutions, independently of
the coefficients. Components with $\min(\alpha,\beta)=1$ are rational curves;
they carry $x^{o(1)}$ solutions by Stewart's bound for Thue equations
\cite{Stewart} (Proposition~\ref{prop:lowdeg}), once a Riemann--Hurwitz argument
shows that the relevant binary form has at least three distinct roots
(Lemma~\ref{lem:threeroots}). Summing over intervals with tail bounds for the
reduced lattices and the Cauchy--Schwarz inequality gives a per-box exponent
$\Bred=m+\tfrac3{10}\lst+\tfrac15\tst$ (Theorem~\ref{thm:Bred}); Reuss's
corresponding exponent is $m+\min(\lst,\tst)$ (Remark~\ref{rem:compare}). The
final exponents are the suprema of explicit piecewise-quadratic functionals,
bounded by exact certificates (Section~\ref{sec:cert}).

For $k\ge3$ the proof uses only three pairwise inputs: the Pell count, the Thue
count and $\Bred$. For $k=2$ it also uses Reuss's Theorem~1, on windows whose
cube-part scales are both at least a fixed power of $x$, and a
Heath-Brown-type bound from \cite[\S5]{Reuss}.

\subsection{What is not claimed}
The exponent $1/4$ of Theorem~\ref{thm:k2} is exactly the supremum of the
functional behind it, attained at a degenerate corner where three bounds that
were known before this paper meet (Remark~\ref{rem:quarter}). We obtain
nothing below $1/4$ for $k=2$. A uniform saving over the exponent $1/10$ for the
components of bidegree $(2,2)$ would lower $77/320$, for example to $89/384$
(Remark~\ref{rem:genus1}); we do not prove such a saving. The constants of
Theorem~\ref{thm:uniform} are not explicit, because the threshold in Tao's
Lemma~3.1 is not; we do not
claim explicit constants anywhere. The formal proofs assume the six results from the
literature stated in Section~\ref{sec:cited}; they do not prove them.

\subsection{Formal status}\label{sec:formalstatus}
All three theorems, and the propositions and lemmas of Sections~\ref{sec:setup}
to~\ref{sec:proofs}, are proved in Lean~4 (version 4.33.1) with Mathlib. The
Lean statements are the theorems above, with the implied constants written as
explicit existential quantifiers, and with the results of
Section~\ref{sec:cited} as named hypotheses (\texttt{let}-bound propositions
assumed in the statement; no axiom is added). Their axiom sets are exactly
\texttt{propext}, \texttt{Classical.choice} and \texttt{Quot.sound}.
Appendix~\ref{app:lean} lists the Lean declaration for each statement, and
Section~\ref{sec:status} describes the checks.

\subsection*{Organisation}
Section~\ref{sec:setup} sets up the pair equations, and Section~\ref{sec:cited}
states the results from the literature with their Lean forms. Sections~\ref{sec:elem}
to~\ref{sec:proofs} give the proof. The steps that carry the new argument
(Lemma~\ref{lem:threeroots}, Proposition~\ref{prop:lowdeg},
Lemma~\ref{lem:P}, Proposition~\ref{prop:interval},
Theorem~\ref{thm:Bred} and Lemma~\ref{lem:goodtriple}) are proved in full,
following the route of the corresponding Lean proof; the remaining steps are
sketched. The Lean declarations listed in Appendix~\ref{app:lean} are the
reference for every statement. Section~\ref{sec:remarks} collects remarks, and
Section~\ref{sec:status} describes verification and the use of AI systems.

\section{Set-up}\label{sec:setup}

\subsection{Notation}
We write $y\sim Y$ for $Y<y\le 2Y$; dyadic scales $Y$ run over
$\{2^j:j\ge-1\}$, so that $y=1$ has the scale $Y=1/2$. Implied constants may
depend on $k$, on $\varepsilon$ and on the auxiliary parameters $\delta,c$
introduced below; any further dependence is indicated. We put
$P_k:=\prod_{p<k}p$. For a polynomial $G$ with integer coefficients, $H(G)$ is
the maximum absolute value of its coefficients.

\subsection{The decomposition}
\begin{lemma}\label{lem:decomp}
Let $k\ge2$ and $n\ge0$ be such that $(n+1)\cdots(n+k)$ is powerful. For
$1\le i\le k$ put
\[
f_i:=\prod_{\substack{p<k\\ v_p(n+i)\ \mathrm{odd}}}p,\qquad
d_i:=\prod_{\substack{p\ge k\\ v_p(n+i)\ \mathrm{odd}}}p .
\]
Then $n+i=f_id_i^3u_i^2$ for a positive integer $u_i$; $f_i\mid P_k$; $d_i$ is
squarefree and coprime to $P_k$. For $i\ne j$ we have $\gcd(d_i,d_j)=1$,
$\gcd(d_j,f_iu_i)=1$ and $\gcd(u_i,u_j)^2\mid (j-i)$.
\end{lemma}

A prime $p\ge k$ divides at most one factor, so its odd valuation there is at
least three; removing $f_id_i^3$ leaves a square.

\begin{lemma}[Pair equation]\label{lem:pair}
In the situation of Lemma~\ref{lem:decomp}, let $1\le i<j\le k$ and put
$d:=f_id_i$, $e:=f_jd_j$, $u:=f_ju_i$, $v:=f_iu_j$ and $h:=f_i^2f_j^2(j-i)$.
Then $e^3v^2-d^3u^2=h$ and $d^3u^2=f_i^2f_j^2(n+i)$. Moreover
$1\le h\le kP_k^4$, $\gcd(d,e)\mid P_k$ and $\gcd(u,v)\le k^{1/2}P_k^2$. For
fixed $(i,j,f_i,f_j)$ the quadruple $(d,e,u,v)$ determines $n$.
\end{lemma}
For $k=2$ this is the set-up of \cite[\S5]{Reuss}, with $h=1$.

\subsection{Windows and exponents}
Fix a nonzero integer $h$. Following \cite[Theorem~1]{Reuss}, for
$D,E\ge1/2$ and $x>1$ put $U:=x^{1/2}D^{-3/2}$, $V:=x^{1/2}E^{-3/2}$,
and call $W=(D,E,x)$ a \emph{window}. Let $\mathcal S_W$ be the set of
$(d,e,u,v)\in\N^4$ with $d\sim D$, $e\sim E$, $U<u<2U$, $V<v<2V$ and
$e^3v^2-d^3u^2=h$. We use the exponents
\begin{equation}\label{eq:exponents}
a=\frac{\log D}{\log x},\quad b=\frac{\log E}{\log x},\quad s=a+b,\quad
t=\frac{\log (UV)}{\log x}=1-\frac{3s}2,\quad
\ub=\frac{1-3a}2,\quad \vb=\frac{1-3b}2,
\end{equation}
so $U=x^{\ub}$ and $V=x^{\vb}$, and
\begin{equation}\label{eq:m}
m=m(s):=\frac98\,st=\frac98\,s\Big(1-\frac{3s}{2}\Big),
\end{equation}
which is concave with maximum $m(1/3)=3/16$.

\begin{lemma}[Covering by windows]\label{lem:windows}
Let $D,E\ge1/2$ and let $x$ be large. Every $(d,e,u,v)\in\N^4$ with
$d\sim D$, $e\sim E$, $x/2<d^3u^2\le 2x$ and $e^3v^2-d^3u^2=h$ lies in
$\mathcal S_{W'}$ for one of $O(1)$ windows $W'=(D',E',x')$ with
$D'\asymp D$, $E'\asymp E$ and $x'\asymp x$. This holds uniformly in shifts
with $|h|\le x/100$.
\end{lemma}
A grid $\{2^{j/8}\}$ for $D',E',x'$ places the ratios $D'^3u^2/x'$ and
$E'^3v^2/x'$ strictly inside $(1,4)$; for $|h|\le x/100$ the same grid works
uniformly in $h$.

\section{Results from the literature}\label{sec:cited}
Six results from the literature enter the Lean statements as named
hypotheses: four from journals and books, and two, Reuss's Theorem~1 and
Tao's Lemma~3.1, from arXiv preprints. For each
we give the statement as used, its source, and the Lean form. For Reuss's
paper we use arXiv:1212.3150v2 and specialise to $k=3$, $l=2$ in his notation.
For Theorems~\ref{thm:reuss1}, \ref{thm:ccdn}, \ref{thm:stewart},
\ref{thm:gelfond} and~\ref{thm:tao} the Lean hypothesis is an instance or a
weakening of the source statement, compared with the source text, so it
holds if the source theorem holds. For Theorem~\ref{thm:hb15} it is an instance of the
statement in Heath-Brown's manuscript of \cite{HBLNM}, applied to $2F$ as
explained after the theorem; the published text has not been compared (Section~\ref{sec:status}). The set counted in
Theorem~\ref{thm:k2} is infinite (the Pell family above); for $k\ge3$ the
counted sets are not known to be nonempty.

\begin{theorem}[{Reuss \cite[Theorem~1]{Reuss}, $k=3$, $l=2$, $h=1$}]\label{thm:reuss1}
Let $h=1$ and let $W=(D,E,x)$ be a window with $D,E>1$ and
$\log(DE)\asymp\log(UV)\asymp\log x$. Define $M>1$ by
$\log M=\frac98\log(DE)\log(UV)/\log x$. Then for $x$ large in terms of
$\varepsilon$, the number of $(d,e,u,v)\in\N^4$ with $D<d<2D$, $E<e<2E$,
$U<u<2U$, $V<v<2V$ and $e^3v^2-d^3u^2=1$ is
\[
\ll_\varepsilon x^{\varepsilon}\min\{(DEM)^{1/2}+D+E,\ (UVM)^{1/2}+U+V\}.
\]
\end{theorem}
Reuss writes $d\sim D$ for $D<d<2D$. The Lean hypothesis fixes the
comparability constants as $\frac1{20}\log x\le\log(DE),\log(UV)\le20\log x$
and requires the bound only on windows with $D,E\ge x^c$, for each fixed
$c>0$, with constants depending on $c$ and $\varepsilon$; so it is an instance
of Theorem~\ref{thm:reuss1}. The declaration \path{reuss_window_bridge}
proves that this open-endpoint form implies the closed-endpoint form
$d\sim D$, $e\sim E$ of $\mathcal S_W$: for fixed $(d,e)$ the equation has at
most one solution with $u\in(U,2U)$, so the endpoints $d=2D$ and $e=2E$ add
$\ll x^{\varepsilon}\min(E,U)+x^\varepsilon\min(D,V)$. Reuss states his
Theorem~1 for every fixed $h\ne0$; his proof begins by reducing to
$(ev,du)=1$, a reduction that is immediate for $h=\pm1$. We use the theorem
only with $h=1$, that is, only for $k=2$. In exponents, Theorem~\ref{thm:reuss1} bounds $\#\mathcal S_W$ by
$x^{R(a,b)+\varepsilon}$ with
\begin{equation}\label{eq:R}
R(a,b):=\min\Big\{\max\Big(\frac{s+m}2,a,b\Big),\ \max\Big(\frac{t+m}2,\ub,\vb\Big)\Big\}.
\end{equation}

\begin{theorem}[{Castryck--Cluckers--Dittmann--Nguyen \cite[Theorem~3]{CCDN}}]\label{thm:ccdn}
For each $N>1$ there is $c(N)$ such that for all $d\ge1$, all integral affine
curves $C\subset\A^N_{\overline{\Q}}$ of degree $d$ and all $B\ge1$,
\[
N_{\mathrm{aff}}(C,B):=\#\{P\in C(\Q)\cap\Z^N:\ |P_i|\le B\}\le c(N)\,d^3B^{1/d}(\log B+d).
\]
\end{theorem}
The Lean hypothesis is the case $N=12$, for every prime ideal
$I\subset\Q[x_1,\dots,x_{12}]$ whose quotient has Krull dimension one; the
degree $d$ is the leading coefficient of the Hilbert polynomial of the
quotient. It therefore covers curves that are irreducible over $\Q$ but not
geometrically irreducible. This is the extension stated in \cite[Corollary~4.1.4
and Remark~4.1.5]{CCDN}: the rational points of such a curve lie on the
intersection with a Galois conjugate, at most $d^2$ points.

\begin{theorem}[{Stewart \cite[p.~796, consequence of Theorem~1]{Stewart}}]\label{thm:stewart}
Let $Q\in\Z[r,s]$ be a binary form of degree $q\ge3$ with content $1$ and
nonzero discriminant, and let $m$ be a nonzero integer. Then
\[
\#\{(r,s)\in\Z^2:\ \gcd(r,s)=1,\ Q(r,s)=m\}\le 2800\,q^{1+\omega(|m|)}.
\]
\end{theorem}
Stewart's Theorem~1 assumes only content~$1$ and nonzero discriminant, so
reducible forms are included, and it counts pairs $(r,s)$, so $(r,s)$ and
$(-r,-s)$ are counted separately; the displayed bound is his consequence of
Theorem~1 with $g=|m|$, the factor $1+1/(8\varepsilon q)$ disappearing as
$\varepsilon\to\infty$ \cite[pp.~795--796]{Stewart}. The Lean hypothesis is
this statement, with ``nonzero discriminant'' expressed as squarefreeness of
the form over $\Q$.

\begin{theorem}[{Heath-Brown \cite[Theorem~15]{HBLNM}, $n=2$; see also \cite[Theorem~2]{ReussThesis}}]\label{thm:hb15}
Let $F\in\Z[x_1,x_2]$ be absolutely irreducible of degree $d\ge2$, let
$\eta>0$ and $B_1,B_2\ge2$. Let $T=\max B_1^{e_1}B_2^{e_2}$, the maximum
over the exponent vectors of the monomials occurring in $F$. Then the integral
zeros of $F$ with $|x_i|\le B_i$ all lie on the zero sets of at most
$\ll_{d,\eta}T^{\eta}\exp(\log B_1\log B_2/\log T)\log(2\|F\|)$
polynomials of degree $O_{d,\eta}(1)$, each coprime to $F$.
\end{theorem}
In \cite[Theorem~15]{HBLNM} (Section~2.6 of the author's manuscript of the
notes), and in its restatement \cite[Theorem~2]{ReussThesis}, the degree is
arbitrary, $B_i\ge1$, the auxiliary polynomials have integer coefficients,
are coprime to $F$ and have degree bounded in terms of $d$ and $\eta$, and the
last factor is $\log\|F\|$, with $\|F\|$ the largest modulus of a coefficient.
Applied to $2F$, which has the same integral zeros, degree, monomials and
irreducibility as $F$, the theorem gives the same conclusion with the factor
$\log\|2F\|=\log(2\|F\|)$. The Lean hypothesis asks only for complex
polynomials not divisible by $F$, implied by integer polynomials coprime to
$F$, and assumes $d\ge2$ and $B_i\ge2$, so it is an instance of the theorem
applied to $2F$.

\begin{theorem}[{Gelfond's inequality, consequence; see \cite[Lemma~1.6.11]{BombieriGubler}}]\label{thm:gelfond}
For each $n$ there is $c_n\ge0$ such that if $G\mid F$ in $\Z[x_1,\dots,x_n]$
with $F\ne0$, then $H(G)\le e^{c_n\deg F}H(F)$.
\end{theorem}
Lemma~1.6.11 of \cite{BombieriGubler} states that if $f=f_1\cdots f_m$ for
complex polynomials in $n$ variables, then
$2^{-d}\prod_j\ell_\infty(f_j)\le\ell_\infty(f)\le2^d\prod_j\ell_\infty(f_j)$,
where $\ell_\infty$ is the largest absolute value of a coefficient and $d$ is
the sum of the partial degrees of $f$. If $F=GK$ in $\Z[x_1,\dots,x_n]$ then
$\ell_\infty(K)\ge1$ and $d\le n\deg F$, which gives the displayed statement
with $c_n=n\log2$; the Lean hypothesis is this statement.

\begin{theorem}[{Tao \cite[Lemma~3.1]{Tao}}]\label{thm:tao}
If $(N+1)\cdots(N+H)$ is powerful, then $H<N$, and
$H\le\exp((\log N)^{2/3+o(1)})$ as $N\to\infty$.
\end{theorem}
The Lean hypothesis is the consequence that for some $N_0$, every
$N\ge N_0$ and $3\le k<N$ with $P(N,k)$ powerful has
$k\le\exp((\log N)^{7/10})$. It is used only in Theorem~\ref{thm:uniform}.
No explicit value of $N_0$ is given in \cite{Tao}.

\section{Elementary pair bounds}\label{sec:elem}

\begin{lemma}\label{lem:pellthue}
Let $h\ne0$ and $G_0\ge1$ be fixed.
\begin{enumerate}[label=\textup{(\roman*)}]
\item For positive integers $d,e$ and $Z\ge2$, the number of $(u,v)\in\N^2$ with
$u,v\le Z$ and $e^3v^2-d^3u^2=h$ is $\ll_\varepsilon(deZ)^{\varepsilon}$.
\item For positive integers $u,v$, the number of $(d,e)\in\N^2$ with
$\gcd(d,e)\le G_0$ and $e^3v^2-d^3u^2=h$ is at most $G_0\cdot2800\cdot3^{1+\omega(|h|)}$.
\end{enumerate}
Consequently $\#\mathcal S_W\ll x^{s+\varepsilon}$ for every window $W$, and
$\ll x^{t}$ if also $\gcd(d,e)\le G_0$ on the counted set.
\end{lemma}
Part (i) is a Pell-type count: the solutions correspond to elements of norm
$he^3$ in a real quadratic order, $\tau(|he^3|)$ ideal classes times $O(\log x)$
units. Part (ii) applies Theorem~\ref{thm:stewart} to the binary cubic
$v^2e'^3-u^2d'^3$ after removing $\gcd(d,e)$.

\begin{lemma}\label{lem:hb}
Let $h=1$ and let $W$ be a window. The number of $(d,e,u,v)\in\mathcal S_W$
with $d$ and $e$ squarefree is
\[
\ll_\varepsilon x^{\varepsilon}\,EU\Big(1+\frac{D}{E^3}\Big)^{1/2}
\quad\text{and}\quad
\ll_\varepsilon x^{\varepsilon}\,DV\Big(1+\frac{E}{D^3}\Big)^{1/2};
\]
in exponents, $H_1(a,b):=b+\ub+\tfrac12\max(0,a-3b)$ and
$H_2(a,b):=a+\vb+\tfrac12\max(0,b-3a)$.
\end{lemma}
This is the estimate displayed after (30) in \cite[\S5]{Reuss} and its mirror
image. Fixing $e$ and $u$, the cube congruence $d^3\equiv-u^{-2}\pmod{e^3}$
leaves $3^{\omega(e)}$ residue classes $c$; writing $d=c+e^3t$, the pairs
$(t,v)$ are integral zeros of the absolutely irreducible cubic
$e^3w^2-u^2(c+e^3t)^3-1$, and Theorem~\ref{thm:hb15} with Bézout's theorem
gives $\ll x^{\varepsilon}(1+D/E^3)^{1/2}$ of them.

\section{The per-interval count}\label{sec:interval}

Throughout this section $h\ne0$ is fixed, $W=(D,E,x)$ is a window with
$s,t\ge c$ for a fixed $c\in(0,1)$, $\delta\in(0,1)$ is fixed, $x$ is large
in terms of $h$, $c$ and $\delta$, and $M:=\lceil x^{(1+\delta)m}\rceil$. Then
$x^c\le DE$, $x^c\le UV$, $x^{c^2}\le M\le x^{1/2}$ and
$\log M\ge\frac98(1+\delta)\log(DE)\log(UV)/\log x$. The range $D/(2E)<\rho\le 2D/E$ of
$\rho=d/e$ is covered by at most $3M+1$ intervals
$I=(\rho_0,\rho_0(1+M^{-1})]$ with
$\rho_0\in\{\frac{D}{2E}(1+M^{-1})^j:j\ge0\}\cap[\frac D{2E},\frac{2D}E]$. We fix a
finite set $\mathcal G$ of positive integers and count solutions with
$\gcd(d,e),\gcd(u,v)\in\mathcal G$; by Lemma~\ref{lem:pair} this covers the
solutions coming from Lemma~\ref{lem:decomp}. Implied constants may depend on
$h$, $c$, $\delta$, $\mathcal G$ and $\varepsilon$.

\subsection{The determinant step}
\begin{lemma}[Reuss's Lemma~7 for fixed $h\ne0$]\label{thm:lemma7}
Let $h\ne0$, $c\in(0,1)$ and $\delta\in(0,1)$ be fixed. There are an integer
$K_0\ge1$ and a real $\kappa>0$, depending only on $c$ and $\delta$, with the
following property. Let $x$ be large in terms of $h$, $c$ and $\delta$, let
$W$ be a window with $DE\ge x^c$ and $UV\ge x^c$, and let $M$ be an
integer with $x^{c^2}\le M\le x$ and
$\log M\ge\frac98(1+\delta)\log(DE)\log(UV)/\log x$. Let
$I=(\rho_0,\rho_0(1+M^{-1})]$ with $D/(2E)\le\rho_0\le2D/E$. Then there
is a nonzero $C_I\in\Z[\rho,\omega]$ with $\deg_\rho C_I,\deg_\omega C_I\le K_0$
and coefficients at most $x^{\kappa}$, such that $C_I(d/e,v/u)=0$ for all
$(d,e,u,v)\in\mathcal S_W$ with $d/e\in I$.
\end{lemma}
This is \cite[Lemma~7]{Reuss}, whose proof begins with the reduction to
$(ev,du)=1$ recalled in Section~\ref{sec:cited}; the reduction is not needed.
The approximation $|v/u-(d/e)^{3/2}|\ll|h|(D/E)^{3/2}/x$ and a Taylor
expansion make every maximal minor of the matrix
$(\rho_j^a\omega_j^b)_{a\le A,b\le B}$ smaller than the reciprocal of its
integral denominators, by Cauchy--Binet and a rearrangement bound for
$\sum(i\ell_1+k\ell_2)$ over $H$ lattice points; the choice $A\log(DE)\approx
B\log(UV)$ makes this work exactly when $\log M$ exceeds
$\frac98\log(DE)\log(UV)/\log x$. A kernel vector of the rank-deficient
integral matrix gives $C_I$, with the given representations $d_j/e_j$, $v_j/u_j$
as denominators, so no coprimality is used.

\subsection{Components of the auxiliary curve}\label{sec:components}
Let $G_I(d,e;v,u):=e^{A}u^{B}C_I(d/e,v/u)$, where $A,B\le K_0$ bound the
degrees of $C_I$. It is a nonzero bihomogeneous form of bidegree $(A,B)$ with
coefficients at most $x^\kappa$, and it vanishes at every solution with
$d/e\in I$. Factor $G_I$ into irreducible elements of $\Z[d,e;v,u]$. There are
at most $A+B\le2K_0$ non-constant irreducible factors; each is bihomogeneous, and by
Theorem~\ref{thm:gelfond} its coefficients are at most
$e^{c_4(A+B)}x^\kappa\le x^{\kappa+1}$. Every solution is a zero of one of
them. By Gauss's lemma a factor $G$ of bidegree $(\alpha,\beta)$ is
irreducible over $\Q$, and it is of one of the following types.
\begin{enumerate}[label=(\alph*)]
\item $\beta=0$ and $\alpha\ge2$, or $\alpha=0$ and $\beta\ge2$. A rational
zero $(w_0,w_1)\ne0$ of a form $G(d,e)$ would give the linear factor
$w_1d-w_0e$; so $G$ has no integral zero with $(d,e)\ne0$ (respectively
$(v,u)\ne0$), and carries no solution.
\item[(b1)] A line $pd+qe=0$ or $pv+qu=0$.
\item[(b2)] $\alpha,\beta\ge1$ and $\min(\alpha,\beta)=1$. Then $G$ is
irreducible over $\C$. Say $\beta=1$ and write $G=vP(d,e)+uQ(d,e)$. A common
factor of $P$ and $Q$ over $\Q$ would split $G$, so $P$ and $Q$ are coprime in
$\Q[d,e]$; a Bézout identity for their dehomogenisations (in either chart)
shows that they have no common zero in $\PP^1(\C)$. If $G=G_1G_2$ over $\C$, bihomogeneity makes one factor,
say $G_1$, of bidegree $(a_1,0)$; if $a_1\ge1$ then $G_1$ vanishes at some
$[r:s]$, so $P(r,s)=Q(r,s)=0$, which is impossible; hence $G_1$ is a
constant. The case $\alpha=1$ is symmetric.
\item[(b3)] $\alpha,\beta\ge2$, so that $2\alpha+3\beta\ge10$.
\end{enumerate}

\subsection{Rational components}\label{sec:lowdeg}

\begin{lemma}\label{lem:threeroots}
Let $\alpha,\beta\ge1$. Let $A,B\in\C[r,s]$ be coprime binary forms of degree
$\beta$ and $C,E\in\C[r,s]$ coprime binary forms of degree $\alpha$, and put
$R:=B^3C^2-A^3E^2$. Suppose $R\ne0$ and that $\beta$ is odd or
$\alpha<3\beta/2$. Then $R$ has at least three distinct zeros in
$\PP^1(\C)$. In particular this holds whenever $\min(\alpha,\beta)=1$.
\end{lemma}

\begin{proof}
Coprime forms of positive degree have no common zero in $\PP^1(\C)$, and none
of $A,B,C,E$ is zero. Suppose that $R$ has at most two distinct zeros. A
linear change of the variables $(r,s)$ preserves all hypotheses, so we may
assume that every zero of $R$ lies in $\{[0:1],[1:0]\}$. Put $a(t)=A(t,1)$,
and $b,c,e$ likewise, so that $r_0(t):=R(t,1)=b^3c^2-a^3e^2$ is a nonzero
polynomial whose only root is $0$; thus $r_0=\kappa t^i$ with $\kappa\ne0$.
For a nonzero polynomial $p$ write $p=t^{o(p)}p_1$ with $p_1(0)\ne0$.

\emph{Claim: $a$ and $b$ are monomials, and
$3o(b)+2o(c)=3o(a)+2o(e)$.} Put $f=b^3c^2$ and $g=a^3e^2$, so $f-g=\kappa
t^i$, and let $W:=tf'-if$. Since $t(\kappa t^i)'=i\kappa t^i$, also
$W=tg'-ig$, and $\deg W\le\min(\deg f,\deg g)$. If $W=0$ then $tf'=if$
forces $f$ to be a multiple of $t^i$, and likewise $g$; then $a,b,c,e$ are
monomials and $o(f)=o(g)=i$. If $W\ne0$, the identity
$W=b^2c\,\bigl(t(3b'c+2bc')-ibc\bigr)$ shows $u_1:=b_1^2c_1\mid W$, and
symmetrically $u_2:=a_1^2e_1\mid W$. Moreover $t^{o(f)}$ and $t^{o(g)}$
divide $W$, hence so does $t^{M_0}$ with $M_0=\max(o(f),o(g))$. A common root
$\tau\ne0$ of $u_1$ and $u_2$ would give $f(\tau)=g(\tau)=0$ and so
$\kappa\tau^i=0$; thus $u_1$, $u_2$ and $t$ are pairwise coprime and
$u_1u_2t^{M_0}\mid W$. Writing $\deg f=o(f)+3\deg b_1+2\deg c_1$ and
$\deg g=o(g)+3\deg a_1+2\deg e_1$, the two inequalities
$\deg u_1+\deg u_2+M_0\le\deg f$ and $\le\deg g$ add up to
$\deg a_1+\deg b_1+(2M_0-o(f)-o(g))\le0$. Hence $a_1,b_1$ are constants and
$o(f)=o(g)$, which is the claim.

\emph{Conclusion.} Let $p=\deg a=o(a)$ and $q=\deg b=o(b)$. As $A,B$ have no
common zero at $[0:1]$, $p=0$ or $q=0$; as they have no common zero at
$[1:0]$, $p=\beta$ or $q=\beta$. Similarly $o(c)=0$ or $o(e)=0$. If $p=0$ and
$q=\beta$, the claim gives $3\beta+2o(c)=2o(e)$, so $o(c)=0$ and
$2o(e)=3\beta$; if $p=\beta$ and $q=0$, it gives $2o(c)=3\beta$. In both
cases $\beta$ is even and $3\beta\le2\alpha$, since $o(c),o(e)\le\alpha$,
contrary to the hypothesis. Finally, if $\min(\alpha,\beta)=1$ then either
$\beta=1$ is odd or $\alpha=1<3\beta/2$.
\end{proof}

\begin{proposition}\label{prop:lowdeg}
Let $\alpha,\beta\ge1$ with $\min(\alpha,\beta)=1$, $K\ge0$ and
$\varepsilon>0$. There is $C$, depending only on $h$, $\mathcal G$, $K$,
$\alpha$, $\beta$ and $\varepsilon$, such that for every $x\ge1$ and every
$G\in\Z[d,e;v,u]$ bihomogeneous of bidegree $(\alpha,\beta)$, irreducible
over $\C$, with coefficients at most $x^K$, the number of $(d,e,u,v)\in\N^4$
with $G(d,e,v,u)=0$, $e^3v^2-d^3u^2=h$ and $\gcd(d,e),\gcd(u,v)\in\mathcal G$
is at most $Cx^{\varepsilon}$. The same holds with $C$ independent of $h$,
uniformly in $0<|h|\le x$.
\end{proposition}
By type (b2) above, this applies to every factor of $G_I$ of that type.

\begin{proof}
Let $\beta=1$. (If $\alpha=1$, write $G=dP(v,u)+eQ(v,u)$, let
$(r,s)=(v,u)/\gcd(u,v)$ and put $A=-Q$, $B=P$, $C=r$, $E=s$; then
$dB(r,s)=eA(r,s)$ and $vE(r,s)=uC(r,s)$, and the argument below applies
verbatim.) Write $G=vP(d,e)+uQ(d,e)$ with forms $P,Q$ of degree $\alpha$,
coprime over $\C$ by irreducibility. Let $(d,e,u,v)$ be counted. Then
$(d,e)\ne0$ because $h\ne0$; let $\gamma=\gcd(d,e)$ and $(r,s)=(d,e)/\gamma$.
Put $A=r$, $B=s$ (as forms of degree $\beta=1$), $C=-Q$ and $E=P$. Then
$dB(r,s)=eA(r,s)$, and dividing $G(d,e,v,u)=0$ by $\gamma^\alpha$ gives
$vE(r,s)=uC(r,s)$.

\emph{Resultants.} For coprime integral forms $P_1,Q_1$ of degree $n\ge1$
with coefficients at most $H$, let $\mathrm{Res}(P_1,Q_1)$ be the product of
the resultants of $P_1(t,1),Q_1(t,1)$ and of $P_1(1,t),Q_1(1,t)$. It is a
nonzero integer, $|\mathrm{Res}|\le((2n)!)^2H^{4n}$ by the Leibniz expansion
of the Sylvester determinants, and for coprime $(r,s)$ the integer
$\gcd(P_1(r,s),Q_1(r,s))$ divides it. Let $g_1=\gcd(A(r,s),B(r,s))$ and
$g_2=\gcd(C(r,s),E(r,s))$; these divide $\mathrm{Res}_1:=\mathrm{Res}(A,B)$
and $\mathrm{Res}_2:=\mathrm{Res}(C,E)$. Write $(A,B)(r,s)=g_1(a',b')$ and
$(C,E)(r,s)=g_2(c',e')$ with coprime pairs. From $dB=eA$ and $vE=uC$ we get
$(d,e)=l(a',b')$ and $(v,u)=\mu(c',e')$ with $|l|=\gcd(d,e)$ and
$|\mu|=\gcd(u,v)$. Substituting into $e^3v^2-d^3u^2=h$ gives
\[
h\,g_1^3g_2^2=R(r,s)\,l^3\mu^2,\qquad R:=B^3C^2-A^3E^2 .
\]
Hence $R(r,s)$ divides the nonzero integer
$M_0:=h\,\mathrm{Res}_1^3\mathrm{Res}_2^2$, and
\[
|M_0|\le|h|\,((2\beta)!)^6((2\alpha)!)^4x^{K(12\beta+8\alpha)}.
\]
In particular
$R\ne0$, so by Lemma~\ref{lem:threeroots} $R$ has at least three distinct
zeros in $\PP^1(\C)$.

\emph{Squarefree kernel.} Let $g\in\Z[t]$ be the primitive polynomial
proportional to the product of the distinct monic irreducible factors of
$R(t,1)$ over $\Q$, and let $Q_R(r,s):=s^{e_0+\deg g}g(r/s)$, where $e_0=1$
if $R$ vanishes at $[1:0]$ and $e_0=0$ otherwise. Then $Q_R$ is an integral
binary form of content $1$, squarefree over $\Q$, of degree $q$ with
$3\le q\le3\beta+2\alpha=:Q_0$ (the lower bound counts the zeros of $R$),
and $Q_R$ divides $R$ in $\Z[r,s]$ by Gauss's lemma. So $Q_R(r,s)$ divides
$M_0$.

\emph{Counting.} The solution is determined by $(r,s)$, $\gcd(d,e)$ and
$\gcd(u,v)$: the pairs $(a',b')$ and $(c',e')$ are determined by $(r,s)$, and
the signs of $l$ and $\mu$ by the positivity of $d,e,u,v$. For each of the
$2d(|M_0|)$ values $m'=\pm k$, $k\mid M_0$, Theorem~\ref{thm:stewart} gives at
most $2800q^{1+\omega(m')}\le2800\,Q_0\,d(|M_0|)^{Q_0}$ coprime pairs with
$Q_R(r,s)=m'$, using $q^{\omega(m')}\le2^{Q_0\omega(M_0)}\le d(|M_0|)^{Q_0}$.
Therefore the count is at most
$5600\,Q_0\,(\#\mathcal G)^2\,d(|M_0|)^{Q_0+1}$, and the divisor bound
$d(n)\ll_\eta n^\eta$ with $|M_0|\ll x^{O(1)}$ gives $Cx^\varepsilon$. The
constant depends on $h$ only through $|h|^\eta$; if $|h|\le x$ this factor is
absorbed into $x^\varepsilon$.
\end{proof}

\subsection{Reduced coordinates}\label{sec:reduced}
\begin{lemma}\label{lem:reduced}
Let $I=(\rho_0,\rho_0(1+M^{-1})]$ with $D/(2E)\le\rho_0\le 2D/E$, and put
$\omega_0=\rho_0^{3/2}$.
\begin{enumerate}[label=\textup{(\roman*)}]
\item If $(d,e,u,v)\in\mathcal S_W$ and $d/e\in I$, then
$0<d-\rho_0e\le 4D/M$ and $|v-\omega_0u|\le17V/M$ (for $x\ge3|h|M$).
\item Let $\Lambda_I=T_I(\Z^2)$ and $\Lambda'_I=T'_I(\Z^2)$, where
$T_I(x_1,x_2)=(\frac{M}{4D}(x_1-\rho_0x_2),\frac{x_2}{2E})$ and
$T'_I(y_1,y_2)=(\frac{M}{17V}(y_1-\omega_0y_2),\frac{y_2}{2U})$. With
Lagrange--Gauss reduced bases $g^{(i)}$, $g'^{(i)}$ put $L_i=2/|g^{(i)}|$ and
$T_i=2/|g'^{(i)}|$. Then $L_1\ge L_2$, $T_1\ge T_2$, $L_1L_2\asymp DE/M$,
$T_1T_2\asymp UV/M$, $L_1\gg(DE/M)^{1/2}$ and $T_1\gg(UV/M)^{1/2}$.
\item Every $(d',e')\in\Z^2$ with $|d'-\rho_0e'|\le4D/M$ and $|e'|\le2E$
has coordinates $\lambda\in\Z^2$, $|\lambda_i|\le L_i$, in the basis
$T_I^{-1}g^{(i)}$ of $\Z^2$, and every $(v',u')$ with
$|v'-\omega_0u'|\le17V/M$ and $|u'|\le2U$ has coordinates $\tau$,
$|\tau_i|\le T_i$, in the basis $T_I'^{-1}g'^{(i)}$. These coordinate changes
lie in $\mathrm{GL}_2(\Z)$.
\end{enumerate}
\end{lemma}
This is Reuss's lattice step \cite[\S2.2]{Reuss} with explicit constants:
(i) follows from $d-\rho_0e\le\rho_0e/M$ and from
$\omega_0(1-1/M)u\le v\le\omega_0(1+3/M)u$, which uses $|h|\le d^3u^2/(3M)$;
(ii) and (iii) follow from the Gram inequalities of a reduced basis, which
give $|\lambda_1g^{(1)}+\lambda_2g^{(2)}|^2\ge\frac12(\lambda_1^2|g^{(1)}|^2+\lambda_2^2|g^{(2)}|^2)$.
In the formal proof the constants are $27DE/M\le L_1L_2\le32DE/M$,
$L_1\ge5(DE/M)^{1/2}$, $117UV/M\le T_1T_2\le136UV/M$ and
$T_1\ge10(UV/M)^{1/2}$.

\subsection{The per-interval CCDN step}\label{sec:P}
\begin{lemma}[Lemma P]\label{lem:P}
Let $F\in\Z[\lambda_1,\lambda_2;\tau_1,\tau_2]$ be bihomogeneous of bidegree
$(3,2)$, let $h\in\Z\setminus\{0\}$, and let $\Gamma\subset\PP^1\times\PP^1$ be
an irreducible curve over $\Q$ of bidegree $(\alpha,\beta)\ne(0,0)$. Let
$L_1\ge L_2>0$ and $T_1\ge T_2>0$. Then the number of
$(\lambda,\tau)\in\Z^2\times\Z^2$ with $\lambda$ and $\tau$ primitive,
$|\lambda_i|\le L_i$, $|\tau_i|\le T_i$, $F(\lambda;\tau)=h$ and
$([\lambda],[\tau])\in\Gamma$ is at most
\[
2\,c(12)\,d^3B^{1/d}(\log B+d),\qquad d:=2\alpha+3\beta,\quad B:=\max(1,L_1^3T_1^2),
\]
with $c(12)$ from Theorem~\ref{thm:ccdn}, uniformly in $F$, $\Gamma$ and $h$.
\end{lemma}

Geometrically, the twelve monomials $\sigma_k$ of bidegree $(3,2)$ embed
$\PP^1\times\PP^1$ by $\mathcal O(3,2)$, and the counted points are integral
points on the affine slice $\{F=h\}$ of the cone over the image of $\Gamma$,
a curve of degree $2\alpha+3\beta$. The proof works with the coordinate ring
of this slice directly.

\begin{proof}
Let $R=\Q[y_0,y_1;y_2,y_3]$, bigraded with $y_0,y_1$ of bidegree $(1,0)$ and
$y_2,y_3$ of bidegree $(0,1)$; let $\sigma_k=y_0^iy_1^{3-i}y_2^jy_3^{2-j}$ for
$k=3i+j$, $0\le i\le3$, $0\le j\le2$, so that $\sigma(\lambda,\tau)\in\Z^{12}$
denotes the vector of these monomials. Let $G$ be an equation of $\Gamma$,
irreducible in $R$, so that $(G)$ is prime. If $G\mid F$, then $F$ vanishes
on $\Gamma$ and, as $h\ne0$, nothing is counted. Otherwise the image $\bar F$
of $F$ in the domain $R/(G)$ is nonzero; let $S=(R/(G))[\bar F^{-1}]$ and
let $\psi\colon\Q[x_1,\dots,x_{12}]\to S$ send $x_k$ to $h\sigma_k/\bar F$.
Its kernel $I$ is a prime ideal.

\emph{Hilbert function.} The image under $\psi$ of the polynomials of total
degree at most $s$ is $\{\bar N\bar F^{-s}:N\in R_{(3s,2s)}\}$. Indeed, a
monomial $x^m$ with $|m|\le s$ maps to
$h^{|m|}\sigma^m\bar F^{\,s-|m|}/\bar F^{\,s}$, whose numerator has bidegree
$(3s,2s)$; conversely every monomial of bidegree $(3s,2s)$ is a product of
$s$ of the $\sigma_k$, and $h\ne0$. Multiplication by $\bar F^{-s}$ is
injective, so this image has dimension $\dim(R/(G))_{(3s,2s)}$. The exact
sequence
$0\to R_{(3s-\alpha,2s-\beta)}\xrightarrow{\ \cdot G\ }R_{(3s,2s)}\to(R/(G))_{(3s,2s)}\to0$
and $\dim R_{(a,b)}=(a+1)(b+1)$ give, for $s\ge\alpha+\beta$,
\[
(3s+1)(2s+1)-(3s-\alpha+1)(2s-\beta+1)=(2\alpha+3\beta)\,s+(\alpha+\beta-\alpha\beta).
\]
So $I$ has degree $d=2\alpha+3\beta$ in the sense of Theorem~\ref{thm:ccdn}.

\emph{Dimension one.} Since the Hilbert function grows linearly, any two
elements of $\Q[x]/I$ are algebraically dependent: the $(k+1)^2$ products
$a^ib^j$ with $i,j\le k$ lie in a filtration piece of dimension $O(k)$. If
$\mathfrak p\ni p\ne0$ is a prime and $\Phi(p,b)=0$ with $\Phi\ne0$, cancelling
the largest power of the first variable dividing $\Phi$ (as $\Q[x]/I$ is a
domain) and reducing modulo $\mathfrak p$ shows that $b$ is algebraic over
$\Q$ modulo $\mathfrak p$. Hence the quotient by any nonzero prime is algebraic
over $\Q$ and therefore a field, so every nonzero prime is maximal; and $\Q[x]/I$, which is
infinite-dimensional and finitely generated, is not a field. So
$\dim\Q[x]/I=1$.

\emph{The counted points.} Let $(\lambda,\tau)$ be counted. Evaluation at
$(\lambda,\tau)$ defines $R/(G)\to\Q$, which sends $\bar F$ to $h\ne0$ and so
extends to $S$; composed with $\psi$ it sends $x_k$ to $\sigma_k(\lambda,\tau)$.
Hence every element of $I$ vanishes at $w=\sigma(\lambda,\tau)\in\Z^{12}$, and
$|w_k|\le L_1^iL_2^{3-i}T_1^jT_2^{2-j}\le L_1^3T_1^2\le B$.

\emph{Fibres.} Suppose $\sigma(\lambda,\tau)=\sigma(\lambda',\tau')$ with all
four vectors primitive. Some $\tau_1^j\tau_2^{2-j}$ is nonzero, and the
coordinates with that $j$ show that the cubic monomial vectors of $\lambda$
and $\lambda'$ are proportional with a nonzero factor; hence $\lambda$ and
$\lambda'$ are proportional, and being primitive, $\lambda'=\kappa\lambda$
with $\kappa=\pm1$. Likewise $\tau'=\mu\tau$ with $\mu=\pm1$. A nonzero
coordinate $\lambda_1^i\lambda_2^{3-i}\tau_1^j\tau_2^{2-j}$ is multiplied by
$\kappa^3\mu^2=\kappa$, so $\kappa=1$. Thus each point $w$ comes from at most
the two pairs $(\lambda,\pm\tau)$.

Therefore the count is at most twice the number of integral points of height
at most $B$ on the curve defined by $I$, and Theorem~\ref{thm:ccdn} with
$N=12$ gives the bound.
\end{proof}

\subsection{The per-interval bound}
\begin{proposition}\label{prop:interval}
In the setting of this section, for each interval $I$ the number of
$(d,e,u,v)\in\mathcal S_W$ with $d/e\in I$ and $\gcd(d,e),\gcd(u,v)\in\mathcal G$ is
\[
\ll_{h,c,\delta,\mathcal G,\varepsilon}x^{\varepsilon}\Big(1+\big(L_1(I)^3T_1(I)^2\big)^{1/10}\Big).
\]
\end{proposition}

\begin{proof}
By Section~\ref{sec:components} every counted solution is a zero of one of at
most $2K_0$ irreducible factors $G$ of $G_I$, each with coefficients at most
$x^{\kappa+1}$. We bound the solutions on each factor. (The formal
statement gives the bound with $L_1$ and $T_1$ computed from some nonzero
vectors of $\Lambda_I$ and $\Lambda'_I$; this implies the form stated, since
the reduced vectors $g^{(1)}$, $g'^{(1)}$ are shortest.)

Type (a) carries none. For a line $pd+qe=0$, $(d,e)$ is $\gamma$ times a fixed
primitive vector with $\gamma\in\mathcal G$, and for each such $(d,e)$
Lemma~\ref{lem:pellthue}(i) gives $\ll x^\varepsilon$ pairs $(u,v)$; for a
line $pv+qu=0$, $(v,u)$ takes at most $\#\mathcal G$ values and
Lemma~\ref{lem:pellthue}(ii) gives $O(1)$ pairs $(d,e)$ for each. Type (b2)
carries $\ll x^\varepsilon$ solutions by Proposition~\ref{prop:lowdeg} with
$K=\kappa+1$.

Let $G$ be of type (b3), and fix $\gamma=\gcd(d,e)$ and
$\gamma'=\gcd(u,v)$ in $\mathcal G$. Put $(d_0,e_0)=(d,e)/\gamma$ and
$(v_0,u_0)=(v,u)/\gamma'$; these are primitive,
$e_0^3v_0^2-d_0^3u_0^2=h':=h/(\gamma^3\gamma'^2)$ is a nonzero integer, and
$G(d_0,e_0,v_0,u_0)=0$ by bihomogeneity. Since $M\le x^{1/2}$ and $x$ is
large, $x\ge3|h|M$, so Lemma~\ref{lem:reduced}(i) gives
$|d_0-\rho_0e_0|\le4D/M$, $|e_0|\le2E$, $|v_0-\omega_0u_0|\le17V/M$ and
$|u_0|\le2U$, and by Lemma~\ref{lem:reduced}(iii) the coordinates $\lambda$,
$\tau$ of $(d_0,e_0)$ and $(v_0,u_0)$ are primitive with $|\lambda_i|\le L_i$
and $|\tau_i|\le T_i$. In these coordinates $e^3v^2-d^3u^2$ becomes a
bihomogeneous form of bidegree $(3,2)$ and $G$ becomes an irreducible form of
the same bidegree, since the changes of variables lie in $\mathrm{GL}_2(\Z)$;
the map to $(\lambda,\tau)$ is injective. Lemma~\ref{lem:P} with $h'$ and
$d=2\alpha+3\beta\ge10$ bounds the solutions by
$2c(12)d^3B^{1/d}(\log B+d)$. Here $B^{1/d}\le1+(L_1^3T_1^2)^{1/10}$, while
$L_1\ll x$ and $T_1\ll x$ give $\log B\ll\log x$, and $d\le5K_0$; so this is
$\ll x^\varepsilon(1+(L_1^3T_1^2)^{1/10})$. Summing over
$(\gamma,\gamma')\in\mathcal G^2$ and over the at most $2K_0$ factors proves
the proposition.
\end{proof}

\section{Summation over intervals}\label{sec:sum}

Put $L_*:=\max((DE/M)^{1/2},32(D+E)/M)$ and $T_*:=\max((UV/M)^{1/2},194(U+V)/M)$.

\begin{lemma}[Tail counts]\label{lem:tails}
For $L\ge32(D+E)/M$, $\#\{I:\ L_1(I)\ge L\}\le720\,DE/L^2$, and for
$T\ge194(U+V)/M$, $\#\{I:\ T_1(I)\ge T\}\le7000\,UV/T^2$.
\end{lemma}
A short lattice vector is a rational approximation to $\rho_0$ or $\omega_0$
with bounded numerator; the points $\rho_0$ are spaced at least $D/(2EM)$
apart, and counting the pairs (denominator, interval) gives the bounds.

\begin{lemma}\label{lem:moments}
For $0<\theta<2$, $\sum_I L_1(I)^{\theta}\ll_\theta M L_*^{\theta}$ and
$\sum_I T_1(I)^{\theta}\ll_\theta M T_*^{\theta}$.
\end{lemma}
The intervals with $L_1(I)\le L_*$ contribute at most $(3M+1)L_*^\theta$, and
by Lemma~\ref{lem:tails} and $DE\le ML_*^2$ the dyadic shell
$2^kL_*\le L_1(I)<2^{k+1}L_*$ contributes $\ll ML_*^\theta(2^\theta/4)^k$; the
series converges because $\theta<2$.

\begin{theorem}[The reduced per-box bound]\label{thm:Bred}
Let $h\ne0$, $c\in(0,1)$ and $\mathcal G$ be fixed, and let $W$ be a window with
$s,t\ge c$. Define
\begin{equation}\label{eq:Bred}
\Bred(a,b):=m+\frac3{10}\,\lst+\frac15\,\tst,
\end{equation}
where $\lst:=\max(\frac{s-m}2,a-m,b-m)$ and $\tst:=\max(\frac{t-m}2,\ub-m,\vb-m)$.
Then for every $\varepsilon>0$ the number of $(d,e,u,v)\in\mathcal S_W$ with
$\gcd(d,e),\gcd(u,v)\in\mathcal G$ is $\ll_{h,c,\mathcal G,\varepsilon}x^{\max(m,\Bred(a,b))+\varepsilon}$.
\end{theorem}

\begin{proof}
Put $\delta=\min(\varepsilon/2,1/2)$ and $M=\lceil x^{(1+\delta)m}\rceil$.
Each solution lies in one of the at most $3M+1$ intervals $I$, so by
Proposition~\ref{prop:interval} (with $\varepsilon/2$) the count is at most
$\sum_I A\,(1+L_1(I)^{3/10}T_1(I)^{1/5})$ with $A\ll x^{\varepsilon/2}$.
For $L,T\ge0$, applying $XY\le X^2+Y^2$ to $X=(L/L_*)^{3/10}$ and
$Y=(T/T_*)^{1/5}$ gives
\[
L^{3/10}T^{1/5}\le L_*^{3/10}T_*^{1/5}\Bigl(\frac{L^{3/5}}{L_*^{3/5}}+\frac{T^{2/5}}{T_*^{2/5}}\Bigr).
\]
Summing over $I$ with Lemma~\ref{lem:moments} for $\theta=3/5$ and
$\theta=2/5$, the count is $\ll A\,(M+ML_*^{3/10}T_*^{1/5})$. Now
$x^{(1+\delta)m}\le M\le2x^{(1+\delta)m}$, $0\le m\le3/16$,
$L_*\ll x^{\lst-\delta m/2}$ and $T_*\ll x^{\tst-\delta m/2}$, so
$M\ll x^{m+\delta}$ and
$ML_*^{3/10}T_*^{1/5}\ll x^{\Bred+\frac34\delta m}\le x^{\Bred+\delta}$. Since
$\delta\le\varepsilon/2$, the count is $\ll x^{\max(m,\Bred)+\varepsilon}$.
\end{proof}

\begin{lemma}\label{lem:Bredm}
If $0\le a,b\le\frac13$ and $s>2/27$, then $\Bred(a,b)\ge m$.
\end{lemma}
\begin{proof}
Here $s\le\frac23$, so $t\ge0$. Then $s-m=s(1-\frac98t)\ge0$ because $s\ge2/27$
is equivalent to $t\le\frac89$, and $t-m=t(1-\frac98s)\ge0$ since $s\le\frac23<\frac89$. So
$\lst,\tst\ge0$.
\end{proof}

\begin{lemma}[Balanced form]\label{lem:balanced}
If $|a-b|\le\frac23m$, then $\Bred(a,b)=\frac34m+\frac1{10}$.
\end{lemma}
\begin{proof}
$a-m\le\frac{s-m}2$ if and only if $a-b\le m$, and $\ub-m\le\frac{t-m}2$ if
and only if $\frac32(b-a)\le m$; with the symmetric conditions this gives
$\lst=\frac{s-m}2$ and $\tst=\frac{t-m}2$, and then
$\Bred=\frac34m+\frac{3s+2t}{20}=\frac34m+\frac1{10}$ because $3s+2t=2$.
\end{proof}
At $a=b=1/6$ one has $s=1/3$, $m=3/16$ and $\Bred=77/320$.

\section{The optimisation and its certificates}\label{sec:cert}

Let the gate $\mathsf G$ be the set of $(a,b)\in[0,1/3]^2$ with $s,t\ge1/10$
and $a,b\ge1/100$, and define
\begin{gather*}
\Phi_3(a,b):=\min\{s,\ t,\ \Bred(a,b)\ \text{if }(a,b)\in\mathsf G\},\\
\Phi_2^\sharp(a,b):=\min\{s,\ t,\ H_1,\ H_2,\ R\ \text{if }(a,b)\in\mathsf G_0,\ \Bred\ \text{if }(a,b)\in\mathsf G\},
\end{gather*}
with $\mathsf G_0=\{(a,b)\in[0,1/3]^2:s,t\ge1/10\}$.

\begin{proposition}\label{prop:certs}
\begin{enumerate}[label=\textup{(\roman*)}]
\item For all $(a_1,a_2,a_3)\in[0,1/3]^3$,
$\min\{\Phi_3(a_1,a_2),\Phi_3(a_1,a_3),\Phi_3(a_2,a_3)\}\le 77/320$, with
equality at $a_1=a_2=a_3=1/6$.
\item For all $(a,b)\in[0,1/3]^2$, $\Phi_2^\sharp(a,b)\le\frac14$, with equality
exactly at $(0,\frac14)$ and $(\frac14,0)$.
\item If $a,a'\in[19/100,21/100]$, then $(a,a')\in\mathsf G$ and
$\Bred(a,a')\le 38059/160000$, with an $O(\eta)$ error on the
$\eta$-neighbourhood for $0\le\eta\le1/1000$.
\end{enumerate}
\end{proposition}
These are finite statements about explicit piecewise-quadratic functions, proved
in Lean by case analysis on the active terms of each maximum and minimum and by
exact polynomial inequalities; part (ii) includes the strip inequality
$\min\{s,H_2\}\le\frac14+\frac a4$ for $a\le1/100$. The same bounds are certified independently by exact-rational branch and
bound (Python \texttt{fractions}) in the supplementary scripts: each box bound
uses the concavity of $m$, $(s+m)/2$ and $(t+m)/2$ in $s$ and the linearity of
the remaining terms. Table~\ref{tab:cert} lists the runs; negative controls
just below each target fail. The point $(a,b)=(1/200,197/800)$
shows that $R$ is needed outside $\mathsf G$: there $\min\{s,H_2\}=201/800>1/4$.

\begin{table}[ht]
\centering
\footnotesize
\begin{tabular}{@{}llllr@{}}
\toprule
Run & functional & region & target $T$ & result (nodes)\\
\midrule
K3\_MIN & $\min_{\mathrm{pairs}}\Phi_3$ & $[0,\frac13]^3$ & $77/320$ & certified (633)\\
K3\_MIN\_CTL & same & same & $0.2406$ & fails near $(\frac16,\frac16,0.171)$\\
K2\_FULL & $\Phi_2^\sharp$ & $[0,\frac13]^2$ & $1/4$ & certified (93)\\
K2\_FULL\_CTL1 & same & same & $\frac14-10^{-6}$ & fails near $(0,\frac14)$\\
CENTRAL & $\min_{\mathrm{pairs}}\Bred$ & $[\frac{19}{100},\frac{21}{100}]^3$ & $38059/160000$ & certified (1)\\
THETA & $\min_{\mathrm{pairs}}\min(s,t,B_{1/12})$ & $[0,\frac13]^3$ & $89/384$ & certified (633)\\
\bottomrule
\end{tabular}
\caption{Exact-rational certificates from \texttt{certificates.py} and
\texttt{endpoint\_cert.py} (K2\_FULL); the outputs, with further runs and
controls, are in the supplementary directory. $B_{1/12}$ is defined in
Remark~\ref{rem:genus1}.}\label{tab:cert}
\end{table}

\begin{proposition}[Central small-gap exclusion]\label{prop:centralgap}
Let $n\ge1$, $H=n+3$, and suppose $(n+1)(n+2)(n+3)$ is powerful. Suppose, for
$r=1,2,3$, that $n+r=f_rd_r^3u_r^2$ with positive integers $u_r$,
$f_r\in\{1,2\}$ and odd squarefree $d_r$, and that $H^{19}<d_r^{100}<H^{21}$. Then
$\max\{|d_1-d_2|,|d_3-d_2|\}>H^{32/1000}$.
\end{proposition}
It is proved in Lean without hypotheses. It excludes only the small-gap part of the central box.

\section{Proofs of the theorems}\label{sec:proofs}

\subsection{Fixed width}
\emph{Theorem~\ref{thm:k3}.} Split $n+1\in(X/2,X]$ dyadically and sort the
$n$ by $(f_1,f_2,f_3)$ and by the dyadic scales $D_i$ of $d_1,d_2,d_3$. For
each of the three pairs, Lemmas~\ref{lem:pair} and~\ref{lem:windows} map the
$n$ of a box injectively into $O(1)$ window sets. With
$\hat a_i=\max(0,\min(\log D_i/\log X,1/3))$, Proposition~\ref{prop:certs}(i)
selects a pair with $s\le77/320$, or $t\le77/320$, or $(\hat a_i,\hat a_j)\in\mathsf G$
and $\Bred\le77/320$; Lemma~\ref{lem:pellthue}, or Theorem~\ref{thm:Bred}
with Lemma~\ref{lem:Bredm}, bounds the box, after a Lipschitz transfer of the exponents by
$O(1/\log X)$. Summing over $O((\log x)^4)$ boxes gives the bound for $N_k$,
and Proposition~\ref{prop:certs}(iii) gives the central subcount.
This proof does not use
Theorems~\ref{thm:reuss1}, \ref{thm:hb15} or~\ref{thm:tao}.

\emph{Theorem~\ref{thm:k2}.} Here $k=2$, $h=1$, and $d=d_1$, $e=d_2$ are
squarefree. On boxes with $\hat a,\hat b\ge\eta$ for a small fixed $\eta>0$,
Proposition~\ref{prop:certs}(ii) selects one of $s,t,H_1,H_2,R,\Bred$ at most
$\frac14$, and Lemmas~\ref{lem:pellthue} and~\ref{lem:hb},
Theorem~\ref{thm:reuss1} (on windows with $D,E\ge x^{c}$) or
Theorem~\ref{thm:Bred} give the bound. On thin boxes $\hat a<\eta$, the strip
inequality $\min\{s,H_2\}\le\frac14+\frac{\hat a}4$ with
Lemmas~\ref{lem:pellthue} and~\ref{lem:hb} gives
$x^{1/4+\eta/4+\varepsilon}$, and $\eta$ is arbitrary; the case $\hat b<\eta$
is symmetric. So Theorem~\ref{thm:reuss1} is needed only on windows with both
cube-part scales at least a fixed power of $x$.

\subsection{Uniformity over all widths}\label{sec:uniform}
Put $\theta=77/320$. Summing Theorem~\ref{thm:k3} over $k$ would not suffice;
the argument instead selects three consecutive factors with small coefficients
and makes the pair bounds uniform in growing shifts.

\begin{lemma}[Three consecutive small kernels]\label{lem:goodtriple}
For $n\ge0$, $k\ge3$ and $1\le j\le k$, let
$f_j=\prod_{p<k,\ v_p(n+j)\text{ odd}}p$. There is $1\le i\le k-2$ with
$f_i,f_{i+1},f_{i+2}\le B(k)$, where $B(k)=30030$ for $k\le14$ and
$B(k)=k^8$ for $k\ge15$.
\end{lemma}

\begin{proof}
If $k\le14$, take $i=1$: every $f_j$ divides $2\cdot3\cdot5\cdot7\cdot11\cdot13=30030$.
Let $k\ge15$. If a prime $p<k$ divides $f_j$, then $p\mid n+j$; there are at
most $\lfloor k/p\rfloor+1\le2\lfloor k/p\rfloor$ such $j$, and
$p^{\lfloor k/p\rfloor}\mid k!$. Hence $\prod_{j=1}^kf_j$ divides $(k!)^2$, so
$\prod_jf_j\le k^{2k}\le k^{8\lfloor k/3\rfloor}$, using $2k\le8\lfloor k/3\rfloor$
for $k\ge15$. If no admissible $i$ existed, each of the $\lfloor k/3\rfloor$
disjoint blocks $\{3l+1,3l+2,3l+3\}$ would contain some $f_j\ge k^8+1$, and
since all $f_j\ge1$ we would get
$(k^8+1)^{\lfloor k/3\rfloor}\le\prod_jf_j\le(k^8)^{\lfloor k/3\rfloor}$, a
contradiction.
\end{proof}

For a window $W$ and an integer $h\ne0$ let $A_W(h)$ be the number of
$(d,e,u,v)\in\mathcal S_W$, defined with shift $h$, such that
$\gcd(d,e)=\gcd(u,v)=1$.

\begin{lemma}[Uniform primitive windows]\label{lem:uniformwindows}
Fix $0<c\le1/10$, $0<\delta\le1/100$ and $\varepsilon>0$. For sufficiently
large $x$, uniformly in the integer shift $1\le |h|\le x^{\delta/4}$ and the
window $W$, the number $A_W(h)$ of solutions in $\mathcal S_W$ with
$\gcd(d,e)=\gcd(u,v)=1$ satisfies $A_W(h)\ll x^{s+\varepsilon}$ and
$A_W(h)\ll x^{t+\varepsilon}$, and if also $s,t\ge c$, then
$A_W(h)\ll_{c,\delta,\varepsilon}x^{\max(m,\Bred(a,b))+\varepsilon+2\delta}$.
\end{lemma}

\begin{proof}
The first two bounds hold for every $x>1$ uniformly in $0<|h|\le x$. Every
solution has $d\le2D$, $e\le2E$ and $u,v\le32x$. For coprime $(d,e)$ the
number of coprime $(u,v)$ with $u,v\le32x$ and $e^3v^2-d^3u^2=h$ is at most
$2\tau(|h|)\bigl(4\log_2((d^3e^3+1)(32x)^2)+2\bigr)$ (the primitive count
behind Lemma~\ref{lem:pellthue}(i)), and for each $(u,v)$
Lemma~\ref{lem:pellthue}(ii) with $G_0=1$ gives at most
$2800\cdot3^{1+\omega(|h|)}$ coprime pairs $(d,e)$. Summing over the at most
$4DE$ pairs $(d,e)$, respectively $4UV$ pairs $(u,v)$, with
$\tau(|h|),3^{\omega(|h|)}\ll_\varepsilon x^{\varepsilon}$, gives the two bounds.
For the third bound, $|h|\le x^{\delta/4}$ controls every place where the
shift enters Sections~\ref{sec:interval} and~\ref{sec:sum}:
it gives $3|h|M\le x$, the hypothesis of Lemma~\ref{lem:reduced}(i); it gives
$|h|\le x$, the range of the second form of Proposition~\ref{prop:lowdeg} and
of the primitive counts of the previous sentences, which replace
Lemma~\ref{lem:pellthue} on the lines of type (b1);
and it gives $3+8|h|\le x^{\delta/3}$, so that in Lemma~\ref{thm:lemma7} the
approximation of $v/u$ by $(d/e)^{3/2}$ loses at most a factor $x^{\delta/3}$
and the determinant estimate holds with degrees and constants depending only
on $c$ and $\delta$. Lemma~\ref{lem:P} does not involve $h$, and since the
solutions are primitive no division by $\gcd(d,e)$ or $\gcd(u,v)$ is needed.
Hence Proposition~\ref{prop:interval} holds with $\mathcal G=\{1\}$ and an
implied constant independent of $h$ in this range, and summing it over the
intervals as in the proof of Theorem~\ref{thm:Bred} gives
$A_W(h)\ll x^{\max(m,\Bred(a,b))+\varepsilon+\delta}$, which implies the
stated bound.
\end{proof}

\begin{lemma}[Three-term count with small coefficients]\label{lem:uniformtriple}
Let $\mathcal Q$ be the positive powerful integers. For every
$\varepsilon>0$ there are $\eta>0$ and $C_\varepsilon$ such that, for
$y\ge2$ and integers $1\le f_1,f_2,f_3\le y^\eta$,
\[
 \#\{0\le N\le y:N+r\in f_r\mathcal Q\ (r=1,2,3)\}
 \le C_\varepsilon y^{\theta+\varepsilon}.
\]
\end{lemma}

\begin{proof}
Take $\delta$ and $\varepsilon'$ small in terms of $\varepsilon$, and $\eta$
small in terms of $\delta$ and $\varepsilon$. The $N\le y^\theta$ contribute at
most $y^\theta+1$. For $y^\theta<N\le y$ write $N+r=f_rd_r^3u_r^2$
($r=1,2,3$) and $F=\max_rf_r\le y^\eta$.

\emph{The primitive equation.} For a pair $(i,j)$ the identity of
Lemma~\ref{lem:pair} gives $e^3v^2-d^3u^2=h$ with $d=f_id_i$, $e=f_jd_j$,
$u=f_ju_i$, $v=f_iu_j$ and $1\le h=f_i^2f_j^2(j-i)\le2F^4$; only this identity
is used, since the $f_r$ are arbitrary. Dividing by $g=\gcd(d,e)$ and
$q=\gcd(u,v)$ gives a primitive solution $(d',e',u',v')$ with shift
$h'=h/(g^3q^2)$, an integer because $g^3q^2\mid h$; so $g,q,h'\le2F^4$, and
its size $d'^3u'^2$ lies between $(N+1)/2$ and $F^4(N+3)$.

\emph{The window and the shift.} For $y$ large,
$h'\le2F^4\le(N+1)/200\le d'^3u'^2/100$, so Lemma~\ref{lem:windows} places the
primitive solution in a window $W'=(D',E',x')$ from the grid $\{2^{k/8}\}$ with
$N+1\ll x'\ll F^4(N+1)$. As $N>y^\theta$ and $F\le y^\eta$, with $\eta$ small
in terms of $\delta$ we get $h'\le2F^4\le x'^{\delta/4}$, so
Lemma~\ref{lem:uniformwindows} applies to $W'$ uniformly in the shift.

\emph{The exponents and the choice of pair.} Let $\hat a_r$ be $\log d_r/\log(N+1)$
clipped to $[0,1/3]$. The exponents $(a,b)$ of $W'$ differ from
$(\hat a_i,\hat a_j)$ by $O((1+\log F)/\log N)$, the clipping included, and
$s,t,m,\lst,\tst$ and $\Bred$ are Lipschitz in $(a,b)$; this perturbation
costs a factor $O(F^{O(1)})$.
Proposition~\ref{prop:certs}(i) at $(\hat a_1,\hat a_2,\hat a_3)$ selects a
pair with $s\le\theta$, or $t\le\theta$, or $(\hat a_i,\hat a_j)\in\mathsf G$
and $\Bred\le\theta$; in the last case $s,t\ge c=1/20$ for $y$ large, and
Lemma~\ref{lem:Bredm} at $(\hat a_i,\hat a_j)$, where $s\ge1/10>2/27$, gives
$m\le\Bred$ there. In each case Lemma~\ref{lem:uniformwindows} bounds the
primitive solutions in $W'$ with shift $h'$ by
$\ll_\varepsilon F^{O(1)}(N+1)^{\theta+\varepsilon'+2\delta}$.

\emph{Summation.} The key consisting of the pair, $g$, $q$ and the window
$W'$ determines $h'$, and together with the primitive solution it determines
$N$. There are $O(F^8(\log y)^3)$ keys. The total is therefore
$y^\theta+1+F^{O(1)}(\log y)^3y^{\theta+\varepsilon'+2\delta}$, and with
$F\le y^\eta$ the choice of parameters makes the exponent less than
$\theta+\varepsilon$.
\end{proof}

\begin{proof}[Proof of Theorem~\ref{thm:uniform}]
Theorem~\ref{thm:tao} applies to every powerful $P(N,k)$ with $N\ge N_0$ and
$k\ge3$: the condition $k<N$ is automatic, because if $N\le k$, Bertrand's
postulate gives a prime $p$ with $N<p\le N+k<2p$, which divides $P(N,k)$
exactly once. So for $n_1=\max(N_0,1)\le n\le x$ every powerful $P(n,k)$ with
$k\ge3$ has $k\le K_x=\exp((\log x)^{7/10})=x^{o(1)}$.

\emph{Width at least three.} Let $x$ be large, so that $14\le K_x\le x$ and
$K_x^8\le x^\eta$. Lemma~\ref{lem:goodtriple} gives an index $i\le K_x$ with
$f_i,f_{i+1},f_{i+2}\le K_x^8$, and $(n+l)/f_l$ is powerful: a prime $p\ge k$
dividing $n+l$ divides no other factor, so $p^2\mid n+l$, and a prime $p<k$
has even exponent in $(n+l)/f_l$. So $N=n+i-1\le2x$ satisfies
$N+r\in f_{i+r-1}\mathcal Q$ for $r=1,2,3$, and Lemma~\ref{lem:uniformtriple} with $\varepsilon/2$ and
$y=2x$ applies to each of the at most $K_x$ indices and $K_x^{24}$ coefficient
triples, and the count is at most
$n_1+K_x^{25}C(2x)^{\theta+\varepsilon/2}\ll_\varepsilon x^{\theta+\varepsilon}$.

\emph{Width two.} If $(n+1)(n+2)$ is powerful, both factors are powerful
since they are coprime, and Theorem~\ref{thm:k2} with $x+1$ applies to $n+1$;
this gives the bound
$x^{1/4+\varepsilon}$ for $k\ge2$.

\emph{Non-powerful members.} A non-powerful $z$ with $N<z\le N+K$ and
$P(N,K)$ powerful forces $K\ge3$, and then $K<N$ by the Bertrand argument.
Starts $N<n_1$ give $z<2n_1$; every other start is counted by the first bound
with $\varepsilon/2$ and carries at most $K\le K_x\le x^{\varepsilon/2}$ such $z$.

The explicit choices of $\delta$, $\varepsilon'$, $\eta$, $c$ and of the
thresholds in $x$ and $y$, and the resulting constants, are those of the
formal proofs of Lemmas~\ref{lem:uniformwindows} and~\ref{lem:uniformtriple}
and Theorem~\ref{thm:uniform} (Appendix~\ref{app:lean}). All implied
constants depend only on $\varepsilon$ and on the threshold $N_0$ of
Theorem~\ref{thm:tao}.
\end{proof}

\section{Remarks}\label{sec:remarks}

\begin{remark}[$1/4$ is the limit of this method for $k=2$]\label{rem:quarter}
At $(a,b)=(0,\frac14)$ every bound used here is at least $\frac14$: there
$m=\frac{45}{256}$, $s=H_2=R=\frac14$, $t=\frac58$, $H_1=\frac34$ and
$\Bred=\frac{673}{2560}$. At this corner $d_1$ is bounded,
$d_2\approx x^{1/4}$, $u_1\approx x^{1/2}$ and $u_2\approx x^{1/8}$. Going
below $\frac14$ at this corner requires a further bound, uniform in $d_1$ and
$u_2$, for the integral points on the genus-one curves $Y^2=d_1u_2^2X^3-d_1$
(a bound $x^{o(1)}$ would suffice); this paper uses none.
\end{remark}

\begin{remark}[Comparison with Reuss]\label{rem:compare}
Reuss's per-box exponent $R(a,b)$ equals $m+\min(\lst,\tst)$. Our exponent is
$m+\frac3{10}\lst+\frac15\tst$. At Reuss's binding point $s=t=2/5$
($m=9/50$, $\lst=\tst=11/100$) these are $0.29$ and $0.235$. The difference
is entirely the per-interval count: Reuss counts rational points of the
auxiliary curve with one coordinate pair at a time \cite[\S2.2, Lemma~8]{Reuss},
whereas Lemma~\ref{lem:P} counts integral points of the curve cut out jointly
by the auxiliary equation and by \eqref{eq:pair}.
\end{remark}

\begin{remark}[A genus-one saving would help]\label{rem:genus1}
If every component of bidegree $(2,2)$ in Proposition~\ref{prop:interval}
carried $\ll x^{\varepsilon}\max(1,L_1^3T_1^2)^{\theta'}$ solutions for some
$\theta'\in[\frac1{12},\frac1{10}]$, the per-box exponent would become
$B_{\theta'}:=m+3\theta'\lst+2\theta'\tst$. For $\theta'=1/12$ the run THETA
certifies $\sup_{[0,1/3]^3}\min_{\mathrm{pairs}}\min(s,t,B_{1/12})=89/384\approx0.2318$.
This would need a uniform saving over the exponent $1/10$ of
Theorem~\ref{thm:ccdn} for integral points on the degree-$10$ curves of
Lemma~\ref{lem:P}. Nothing in this remark is claimed; the value $89/384$ is an
exact-rational certificate, not a Lean theorem.
\end{remark}

\begin{remark}[Relation to Tao's model problem]\label{rem:tao}
Tao \cite[Remark~2.12]{Tao} proposes to show that the variety
$n_1^2n_2^3+1=m_1^2m_2^3$ \cite[(2.9)]{Tao} has $o(H^2)$ integral points of
height at most $H$. The pair equation \eqref{eq:pair} is this variety with
$h=1$, and the geometry of Lemma~\ref{lem:P} is related to it; we have not
examined what our method gives for the model problem as posed, where the cube
variables need not be squarefree. Theorem~\ref{thm:uniform} improves the
quantity in Tao's Theorem~1.8 through the uniform triple count instead.
\end{remark}

\begin{remark}[The constant $9/8$]\label{rem:98}
The exponents depend on the determinant constant through $m=c_0st$ with
$c_0=9/8$; in the balanced case at $s=1/3$ the bound is $c_0/8+1/10$. Within
the determinant estimate of Lemma~\ref{thm:lemma7} (a rectangular monomial set
and a generalised Vandermonde bound), the AM--GM inequality shows that no
choice of $A$ and $B$ lowers $9/8$.
This remark is a written argument; it is not formalized.
\end{remark}

\section{Verification and disclosure}\label{sec:status}

\subsection*{Trust base}
The formal proofs are checked by the kernel of Lean~4.33.1 with Mathlib at
commit \path{0df444a360eaa60ab8c11dca51a86af692955474}. Every declaration
of Table~\ref{tab:finals} depends exactly on the axioms \texttt{propext},
\texttt{Classical.choice} and \texttt{Quot.sound}; there is no
\texttt{sorry}, no \texttt{native\_decide} and no added axiom. The results
from the literature enter only as the named hypotheses of Section~\ref{sec:cited}. What
must be trusted beyond the kernel is therefore: those six results;
the agreement of each Lean hypothesis with its source, discussed in
Section~\ref{sec:cited}; and the agreement of the Lean statements of
Table~\ref{tab:finals} with Theorems~\ref{thm:k2}--\ref{thm:uniform}.

\subsection*{Statements and independent checks}
Each statement in Table~\ref{tab:finals} is stated independently of its
proof, and the proofs were checked against it by the Comparator tool, which
compares the compiled statement with an independently compiled copy, together
with an export of the proof checked by the independent kernel Nanoda with only
the three standard axioms permitted. The declarations are in the module
\path{CountingClosure} of the accompanying repository, compiled with zero
warnings and audited for their exact axiom sets. The replay recorded in
\path{replay/receipts/powerful-counting.json} checks every release
declaration of this paper, the nine of Table~\ref{tab:finals} and the
supporting declarations of Table~\ref{tab:concordance}, with Nanoda (standard
axioms only), and compares each against its independent statement file in
\path{replay/statements/powerful-counting/}; a negative control, the statement
of Theorem~\ref{thm:k3} with $77/320$ replaced by $76/320$, is rejected. The Lean names of
intermediate steps locate their proofs; only the declarations of
Table~\ref{tab:finals} are compared with stated paper results.

\subsection*{Hypotheses from the literature}
The statements of Theorems~\ref{thm:reuss1}, \ref{thm:ccdn} and \ref{thm:tao}
were compared with the arXiv texts \cite{Reuss,CCDN,Tao}, and
Theorem~\ref{thm:stewart} with the journal text \cite[pp.~795--796]{Stewart}
and Theorem~\ref{thm:gelfond} with \cite[Lemma~1.6.11]{BombieriGubler}, by an
AI reviewer. Theorem~\ref{thm:hb15} was compared with Theorem~15 in the
author's manuscript of \cite{HBLNM} and with its restatement in
\cite[Theorem~2]{ReussThesis}; the published text of \cite{HBLNM} has not
been compared.

\subsection*{Use of AI systems}
All results, proofs (written and formal), computation scripts and the text of
this manuscript were produced by AI systems. The results and their proofs come
from large language models running as agents in the Jocelyn research harness,
together with Lean and exact-arithmetic tools; much of the writing and revision
of the manuscript, its reviews and the resulting corrections were done with AI
coding assistants from several providers, operated by the author outside the
harness. The author chose the problem, operated the harness, set its
objectives, decided the content of the manuscript and commissioned the
reviews, and takes responsibility for the manuscript. All reviews were carried
out by AI systems; no human has reviewed the results.

\subsection*{Status}
The novelty of the results has been checked only by bounded literature
searches (arXiv, zbMATH in part, the papers citing Reuss, and Tao's paper and
its references), not certified.

\subsection*{Reproduction}
The Lean sources, statement files and replay receipt are in release
\texttt{v1.0.0} of the accompanying repository,
\url{https://github.com/dannyply/math-frontier}; the receipt
\path{replay/receipts/powerful-counting.json} records the commit at which it
was produced. There \path{scripts/build.sh} compiles
the library and fails on any warning, \path{scripts/audit.sh} checks the axioms
and the statement hash of every release declaration against
\path{release/manifest.json}, and \path{scripts/replay.sh} runs the
independent replay. The certificate scripts are reproduced by
\texttt{python3 certificates.py \&\& python3 endpoint\_cert.py} in the
supplementary directory, in about a minute on a laptop; all arithmetic is
exact.

\section*{Acknowledgements}
The author thanks the Mathlib community for the formal library used in the
accompanying repository.

\appendix
\section{Lean declarations}\label{app:lean}
Table~\ref{tab:finals} lists the release declarations of the counting closure;
each has an independent statement file and passed the replay of
Section~\ref{sec:status}. Table~\ref{tab:concordance} gives the Lean
declaration behind every statement of the paper. Names are relative to
\path{JocelynWorkspace.Formalization.PowerfulCounting}, except those beginning
with \texttt{Erdos137}, which are relative to \path{JocelynWorkspace};
\texttt{CC} is \texttt{CountingClosure} and \texttt{CC.Final} is
\texttt{CountingClosure.Erdos137Final}. A module name followed by eight
hexadecimal digits denotes its declaration \texttt{JocelynClaims.claim\_}$\langle$hash$\rangle$
whose hash begins with those digits; the release declarations are listed in
full in \path{release/manifest.json} of the repository, and a hash prefix
locates any other declaration in the Lean sources. Intermediate
declarations are checked as part of the dependency closure of the release
declarations; only the release declarations are compared with stated results.

\begin{table}[ht]
\centering\footnotesize
\begin{tabular}{@{}l>{\raggedright}p{5.6cm}>{\raggedright\arraybackslash}p{4.4cm}@{}}
\toprule
Result & Lean declaration & Named hypotheses\\
\midrule
Theorem~\ref{thm:k3} & \path{counting_thm_k3} & Stewart, CCDN, Gelfond\\
Theorem~\ref{thm:k2} & \path{counting_thm_k2_reuss} & Stewart, CCDN, Gelfond, Reuss, Heath-Brown\\
Theorem~\ref{thm:uniform} & \path{counting_thm_uniform_reuss} & the five above and Tao\\
Proposition~\ref{prop:interval} & \path{counting_prop_interval} & Stewart, CCDN, Gelfond\\
Theorem~\ref{thm:Bred} & \path{counting_thm_Bred} & Stewart, CCDN, Gelfond\\
Proposition~\ref{prop:interval}, uniform in $h$ & \path{counting_UPI} & Stewart, CCDN, Gelfond\\
Lemma~\ref{lem:uniformwindows}, third bound & \path{counting_lem_uniformwindows_TB} & Stewart, CCDN, Gelfond\\
Lemma~\ref{lem:uniformtriple} & \path{counting_lem_uniformtriple} & Stewart, CCDN, Gelfond\\
Reuss window convention & \path{reuss_window_bridge} & none\\
\bottomrule
\end{tabular}
\caption{Release declarations of the counting closure.}\label{tab:finals}
\end{table}

\begin{table}[ht]
\centering\footnotesize
\begin{tabular}{@{}>{\raggedright}p{3.4cm}>{\raggedright\arraybackslash}p{11.4cm}@{}}
\toprule
Statement & Lean declarations\\
\midrule
Theorem~\ref{thm:k2} & \texttt{CC.Final.counting\_thm\_k2\_reuss}; \texttt{CC.C12} \texttt{a9994a76}, \texttt{CC} \texttt{6ee9f26a}\\
Theorem~\ref{thm:k3} & \texttt{CC.Final.counting\_thm\_k3}; \texttt{CC.C11} \texttt{d974288e}\\
Theorem~\ref{thm:uniform} & \texttt{CC.Final.counting\_thm\_uniform\_reuss}; \texttt{CC} \texttt{3103e62f}, \texttt{e09c0d72}, \texttt{CC.C23} \texttt{8e090865}, \texttt{AllWidthFromUniformTriple} \texttt{2be59d9f}, \texttt{AllWidthFromPairAndUnion} \texttt{fdb026ef}, \texttt{NonpowerfulMembers} \texttt{ceda7068}\\
Lemmas~\ref{lem:decomp}, \ref{lem:pair} & \texttt{DecompositionAndEnvelope} \texttt{9673eb5d}\\
Lemma~\ref{lem:windows} & \texttt{StrictWindows} \texttt{f0cfdddd} (fixed $h$); growing shifts: \texttt{CC.C21.PcPellWin.growing\_cover}\\
Reuss's window convention & \texttt{CC.Final.reuss\_window\_bridge}\\
Lemma~\ref{lem:pellthue} & \texttt{CC.C08} \texttt{ca19e5d2}\\
Lemma~\ref{lem:hb} & \texttt{CC.C13} \texttt{a92c7c0a}\\
Lemma~\ref{thm:lemma7} & \texttt{Paper3OriginalActualAuxiliaryPolynomial} \texttt{1d348eff}\\
Section~\ref{sec:components} & \texttt{CC.C02} \texttt{7327564c}\\
Lemma~\ref{lem:threeroots} & \texttt{CC.C01.PcCurves.three\_roots\_gen}\\
Proposition~\ref{prop:lowdeg} & \texttt{CC.C01} \texttt{6be17294}; uniform in $h$: \texttt{CC.C16} \texttt{b07245cc}\\
Lemma~\ref{lem:reduced} & \texttt{CC.C09} \texttt{cc873556}, \texttt{CC.C04} \texttt{f5611586}\\
Lemma~\ref{lem:P} & \texttt{CC.C03} \texttt{3ede4992}; type (b3): \texttt{CC.C03} \texttt{f159a569}\\
Proposition~\ref{prop:interval} & \texttt{CC.Final.counting\_prop\_interval}; \texttt{CC.C05} \texttt{d9675218}, \texttt{CC.C06} \texttt{40f41d3d}, \texttt{CC.C07} \texttt{71200785}\\
Lemmas~\ref{lem:tails}, \ref{lem:moments} & \texttt{LatticeTails} \texttt{d31e9b77}, \texttt{LatticeMoments} \texttt{de91d4e3}\\
Theorem~\ref{thm:Bred} & \texttt{CC.Final.counting\_thm\_Bred}; \texttt{CC.C10} \texttt{664a2185}, \texttt{8d943e1c}\\
Lemma~\ref{lem:Bredm} & \texttt{DecompositionAndEnvelope} \texttt{928cf894}\\
Lemma~\ref{lem:balanced} & \texttt{DecompositionAndEnvelope.JocelynExpAlg.balanced\_aux}\\
Proposition~\ref{prop:certs} & (i) \texttt{TripleExponentCertificate} \texttt{d17527ab}; (ii) \texttt{PairExponentCertificate} \texttt{c433f729}; (iii) \texttt{DecompositionAndEnvelope} \texttt{9673eb5d}, \texttt{f1b2bf39}\\
Proposition~\ref{prop:centralgap} & \texttt{Erdos137.WidthThree.CentralSet.SmallGapExclusion.small\_gap\_exclusion}\\
Lemma~\ref{lem:goodtriple} & \texttt{AllWidthFromUniformTriple.JocelynAllWidth.goodtriple} (same statement: \texttt{GoodTriple} \texttt{6ba68896}); $k\ge15$: \texttt{Erdos137.FixedWidth.GoodIndex.three\_consecutive\_small\_kernels}\\
Lemma~\ref{lem:uniformwindows} & \texttt{CC.C08.PcPell.uniform\_s}, \texttt{CC.C08.PcPell.uniform\_t} (first two bounds); \texttt{CC.Final.counting\_lem\_uniformwindows\_TB} (third bound), via \texttt{CC.C19} \texttt{62fb44f0} (\texttt{counting\_UPI}), \texttt{CC.C10} \texttt{8d943e1c}, \texttt{CC.C20} \texttt{88a200ce}\\
Lemma~\ref{lem:uniformtriple} & \texttt{CC.Final.counting\_lem\_uniformtriple}; \texttt{CC.C21} \texttt{b14c9f52}\\
\bottomrule
\end{tabular}
\caption{Concordance between the statements of the paper and the Lean development.}\label{tab:concordance}
\end{table}

\begin{thebibliography}{99}

\bibitem{AktasMurty} K.~Akta\c{s} and M.~R.~Murty, \emph{Fundamental units and
consecutive squarefull numbers}, Int.~J. Number Theory \textbf{13} (2017),
243--252.

\bibitem{Bloom} T.~F.~Bloom, \emph{Erd\H{o}s problems},
\url{https://www.erdosproblems.com}, Problems \#137, \#364, \#365.

\bibitem{BombieriGubler} E.~Bombieri and W.~Gubler, \emph{Heights in
Diophantine Geometry}, Cambridge University Press, 2006.

\bibitem{CCDN} W.~Castryck, R.~Cluckers, P.~Dittmann and K.~H.~Nguyen,
\emph{The dimension growth conjecture, polynomial in the degree and without
logarithmic factors}, Algebra Number Theory \textbf{14} (2020), 2261--2294;
arXiv:1904.13109.

\bibitem{Chan2012} T.~H.~Chan, \emph{Twin squareful numbers}, J.~Aust. Math.
Soc. \textbf{93} (2012), 43--51.

\bibitem{ErdosManitoba} P.~Erd\H{o}s, \emph{Problems and results on number
theoretic properties of consecutive integers and related questions}, in:
Proc. Fifth Manitoba Conf. Numer. Math. (1975), 1976, 25--44.

\bibitem{Erdos76} P.~Erd\H{o}s, \emph{Problems and results on consecutive
integers}, Publ. Math. Debrecen \textbf{23} (1976), 272--282.

\bibitem{ErdosGraham} P.~Erd\H{o}s and R.~L.~Graham, \emph{Old and New Problems
and Results in Combinatorial Number Theory}, Monographies de L'Enseignement
Mathématique \textbf{28}, Geneva, 1980.

\bibitem{ErdosSelfridge} P.~Erd\H{o}s and J.~L.~Selfridge, \emph{The product of
consecutive integers is never a power}, Illinois J. Math. \textbf{19} (1975),
292--301.

\bibitem{Golomb} S.~W.~Golomb, \emph{Powerful numbers}, Amer. Math. Monthly
\textbf{77} (1970), 848--852.

\bibitem{HBAnnals} D.~R.~Heath-Brown, \emph{The density of rational points on
curves and surfaces}, Ann. of Math. (2) \textbf{155} (2002), 553--595.

\bibitem{HBLNM} D.~R.~Heath-Brown, \emph{Counting rational points on algebraic
varieties}, in: Analytic Number Theory (C.I.M.E. Summer School, Cetraro, 2002),
Lecture Notes in Math. \textbf{1891}, Springer, 2006, 51--95; author's
manuscript: Oxford Mathematical Institute eprint 676.

\bibitem{HBsqfree} D.~R.~Heath-Brown, \emph{Square-free values of $n^2+1$},
Acta Arith. \textbf{155} (2012), 1--13; arXiv:1010.6217.

\bibitem{Reuss} Thomas Reuss, \emph{Pairs of $k$-free numbers, consecutive
square-full numbers}, arXiv:1212.3150v2 (2014).

\bibitem{ReussThesis} Thomas Reuss, \emph{The Determinant Method and
Applications}, DPhil thesis, University of Oxford, 2015.

\bibitem{Stewart} C.~L.~Stewart, \emph{On the number of solutions of
polynomial congruences and Thue equations}, J.~Amer. Math. Soc. \textbf{4}
(1991), 793--835.

\bibitem{Tao} T.~Tao, \emph{Products of consecutive integers with unusual
anatomy}, arXiv:2603.27990v3 (2026).

\bibitem{Walker} D.~T.~Walker, \emph{Consecutive integer pairs of powerful
numbers and related Diophantine equations}, Fibonacci Quart. \textbf{14}
(1976), 111--116.

\bibitem{WongcharoenbhornMeemark} W.~Wongcharoenbhorn and Y.~Meemark,
\emph{Square-full values of quadratic polynomials}, Bull. Aust. Math. Soc.
\textbf{113} (2026), 40--54; arXiv:2405.06968.

\end{thebibliography}

\end{document}
